\documentclass[12pt,a4paper]{article}

% ----------- Encoding / language -----------
\usepackage[utf8]{inputenc}
\usepackage[T1]{fontenc}
\usepackage[english]{babel}

% ----------- Math / layout -----------
\usepackage{amsmath,amssymb}
\usepackage{geometry}
\usepackage{setspace}
\usepackage{titlesec}

% ----------- Figures / floats -----------
\usepackage{graphicx}
\usepackage{float}
\usepackage{caption}

% ----------- Header/footer -----------
\usepackage{fancyhdr}

% ----------- Lists -----------
\usepackage{enumitem}

% ----------- Fonts -----------
\usepackage{newtxtext,newtxmath}

% ----------- Code listings -----------
\usepackage{listings}
\usepackage{courier}
\usepackage{xcolor}

% ----------- Tables -----------
\usepackage{booktabs}
\usepackage{array}

% =============== Page setup ===============
\geometry{top=2cm,bottom=2cm,left=2.5cm,right=2.5cm}
\setstretch{1.5}

\titleformat{\section}{\normalfont\Large\bfseries}{\thesection}{1em}{}
\titleformat{\subsection}{\normalfont\large\bfseries}{\thesubsection}{1em}{}

\pagestyle{fancy}
\fancyhf{}
\rfoot{\thepage}

\captionsetup[figure]{position=bottom,justification=centering}

\definecolor{codegray}{rgb}{0.95,0.95,0.95}
\definecolor{commentcolor}{rgb}{0.4,0.4,0.4}
\definecolor{keywordcolor}{rgb}{0.1,0.1,0.6}
\definecolor{stringcolor}{rgb}{0.2,0.5,0.2}

\lstset{
  language=[x86masm]Assembler,
  basicstyle=\ttfamily\small,
  backgroundcolor=\color{codegray},
  commentstyle=\color{commentcolor}\itshape,
  keywordstyle=\color{keywordcolor}\bfseries,
  stringstyle=\color{stringcolor},
  numbers=left,
  numberstyle=\tiny\color{gray},
  stepnumber=1,
  numbersep=8pt,
  frame=single,
  rulecolor=\color{black!30},
  showstringspaces=false,
  tabsize=4,
  breaklines=true,
  captionpos=b
}

\graphicspath{{images/}}

\newcommand{\img}[3][0.8]{%
  \begin{figure}[H]
    \centering
    \IfFileExists{#2}{%
      \includegraphics[width=#1\textwidth]{#2}%
    }{%
      \fbox{\parbox{0.85\textwidth}{\centering\textbf{Image missing:}
      \texttt{#2}\\Place the file in the \texttt{images/} folder.}}%
    }%
    \caption{#3}
  \end{figure}
}

\newcommand{\code}[1]{\texttt{\small #1}}

\begin{document}

% ===================== TITLE PAGE =====================
\begin{titlepage}
\begin{center}
  \IfFileExists{utm_logo_1.jpg}{%
    \includegraphics[width=0.23\textwidth]{utm_logo_1.jpg}\\[0.8em]
  }{}

  \textbf{MINISTRY OF EDUCATION AND RESEARCH}\\
  \textbf{OF THE REPUBLIC OF MOLDOVA}\\[0.5em]
  \textbf{TECHNICAL UNIVERSITY OF MOLDOVA}\\
  \textbf{Faculty of Computers, Informatics and Microelectronics}\\
  \textbf{Department of Software Engineering and Automatics}\\[2em]

  \vspace{1.5cm}
  \textbf{\Large Tacu Igor \textemdash{} FAF-242}\\[3em]

  \textbf{\Large REPORT}\\[0.5em]
  Laboratory Work No.\ 5\\[0.5em]
  \textbf{Computer Architecture / Assembly Language}\\[0.5em]
  \textbf{NASM --- Arithmetic Expressions and BCD Numbers}\\[3em]

  \vfill
  \begin{flushright}
    Lecturer:\\
    asist.\ univ.\ Olga Grosu\\
    DISA, FCIM, UTM
  \end{flushright}
  \vfill

  \large \textbf{Chi\c{s}in\u{a}u \textemdash{} 2025}
\end{center}
\end{titlepage}

\tableofcontents
\newpage

% ===================== 1. INTRODUCTION =====================
\section{Introduction}

\subsection{Overview of Lab Topics}

This laboratory work extends the Assembly programming skills developed in Lab 4
by focusing on two new areas: \textbf{signed integer arithmetic expressions} in
x86-64 Assembly, and \textbf{BCD (Binary Coded Decimal)} number representation
and arithmetic.

Signed arithmetic in Assembly requires understanding how two's complement
representation works in hardware. The x86-64 ISA provides instructions such as
\code{IMUL} (signed multiply), \code{IDIV} (signed divide), \code{NEG} (negate),
and \code{CQO} (sign-extend RAX into RDX:RAX for division setup). Correct use
of these requires careful attention to register state and overflow conditions.

BCD encoding is an alternative to pure binary: each decimal digit (0--9) is
individually encoded in four bits. Packed BCD fits two digits into one byte
(one per nibble). While modern CPUs use binary arithmetic internally, BCD
remains relevant in financial software, embedded controllers, and display
drivers where preserving exact decimal representations is required.

\subsection{Variant Selection (Student Number 25)}

The lab specifies a rule to select three expression variants from the list of
fifteen. Starting from student number 25:

\begin{enumerate}
  \item $25 > 15 \Rightarrow 25 / 2 = 12.5 \Rightarrow$ round to \textbf{13}
  \item $13$ is odd $\Rightarrow 13 + 3 = 16 > 15 \Rightarrow 16 / 2 = $ \textbf{8}
  \item $8$ is even $\Rightarrow 8 + 2 = $ \textbf{10}
\end{enumerate}

The three selected variants are \textbf{13}, \textbf{8}, and \textbf{10}.

\newpage

% ===================== 2. PURPOSE =====================
\section{Purpose of the Report}

This report documents the following lab tasks:

\begin{enumerate}
  \item Write an Assembly program computing $A = (A + B) - (C + D)$ using
        registers, with full comments.
  \item Write an Assembly program evaluating three arithmetic expressions
        (variants 13, 8, 10), with register annotations and verification.
  \item Perform manual BCD conversions and additions for the selected variants.
  \item Describe the design of a BCD adder circuit in Logisim.
  \item Write an Assembly program that adds two 8-bit packed BCD numbers and
        stores the result (including carry) in subsequent memory locations.
\end{enumerate}

\newpage

% ===================== 3. ASSEMBLY BASICS REVIEW =====================
\section{Relevant Assembly Concepts}

\subsection{Signed Arithmetic Instructions}

\begin{table}[H]
  \centering
  \caption{Key x86-64 signed arithmetic instructions}
  \renewcommand{\arraystretch}{1.3}
  \begin{tabular}{|l|l|}
    \hline
    \textbf{Instruction} & \textbf{Effect} \\
    \hline
    \code{NEG rax}           & Two's complement negation: \code{rax = -rax} \\
    \code{IMUL rax, rbx}     & Signed multiply: \code{rax = rax * rbx} \\
    \code{CQO}               & Sign-extend RAX into RDX:RAX (required before IDIV) \\
    \code{IDIV rbx}          & Signed divide: \code{rax = rdx:rax / rbx}, remainder in RDX \\
    \code{ADD rax, rbx}      & \code{rax = rax + rbx} (same for signed and unsigned) \\
    \code{SUB rax, rbx}      & \code{rax = rax - rbx} (same for signed and unsigned) \\
    \hline
  \end{tabular}
\end{table}

\noindent\textbf{Why CQO before IDIV?} The \code{IDIV} instruction divides the
128-bit value \code{RDX:RAX} by the operand. If RDX is not properly
sign-extended from RAX, the high bits contain garbage, causing a \emph{divide
error} (\#DE) fault. \code{CQO} (Convert Quadword to Octword) propagates the
sign bit of RAX into every bit of RDX, making the pair a valid signed 128-bit
value.

\subsection{BCD Number Representation}

\textbf{Binary Coded Decimal (BCD)} encodes each decimal digit as a 4-bit
binary group. The 8421 code is the standard variant:

\begin{table}[H]
  \centering
  \caption{8421 BCD code for digits 0--9}
  \renewcommand{\arraystretch}{1.2}
  \begin{tabular}{|c|c|}
    \hline
    \textbf{Decimal digit} & \textbf{8421 BCD (4 bits)} \\
    \hline
    0 & 0000 \\
    1 & 0001 \\
    2 & 0010 \\
    3 & 0011 \\
    4 & 0100 \\
    5 & 0101 \\
    6 & 0110 \\
    7 & 0111 \\
    8 & 1000 \\
    9 & 1001 \\
    \hline
    10--15 & \textit{invalid states} \\
    \hline
  \end{tabular}
\end{table}

\noindent\textbf{Packed BCD} stores two digits per byte using the high nibble
for the tens digit and the low nibble for the units digit. For example,
decimal $57 = \texttt{0x57}$ in packed BCD because
$5 = \texttt{0101}_2$ and $7 = \texttt{0111}_2$.

\subsection{BCD vs.\ Pure Binary: Bit Comparison}

\begin{table}[H]
  \centering
  \caption{Bit-width comparison: binary vs.\ packed BCD for lab variants}
  \renewcommand{\arraystretch}{1.3}
  \begin{tabular}{|c|c|c|c|c|}
    \hline
    \textbf{Decimal} & \textbf{Binary} & \textbf{Bits (bin)} & \textbf{Packed BCD} & \textbf{Bits (BCD)} \\
    \hline
    8  & \texttt{1000}          & 4 & \texttt{0000\ 1000} & 8 \\
    10 & \texttt{1010}          & 4 & \texttt{0001\ 0000} & 8 \\
    13 & \texttt{1101}          & 4 & \texttt{0001\ 0011} & 8 \\
    \hline
  \end{tabular}
\end{table}

\noindent BCD uses exactly twice as many bits as binary for single-digit
numbers and about 20\% more bits overall for large numbers.  The tradeoff is
exact decimal representation with simpler conversion to display form.

\subsection{BCD Addition Rules}

\begin{enumerate}
  \item Add the two BCD numbers using ordinary binary addition.
  \item For each 4-bit group (nibble):
    \begin{itemize}
      \item If the nibble result is $\leq 9$ and no carry was generated out
            of that nibble, it is a valid BCD digit --- no adjustment needed.
      \item If the nibble result is $> 9$ (i.e., in the invalid range
            \texttt{0xA}--\texttt{0xF}), \textbf{add 6} to skip the six
            invalid states and wrap back into the valid BCD range; the +6
            produces a carry into the next nibble.
      \item If a carry \emph{was} generated out of the nibble (even if the
            result was $\leq 9$), \textbf{add 6} as well.
    \end{itemize}
\end{enumerate}

\newpage

% ===================== 4. TASK 1 =====================
\section{Task 1: Expression $A = (A+B)-(C+D)$}

\subsection{Specification}

Write an Assembly program that computes $A = (A + B) - (C + D)$ using
registers.  Values used: $A = 10,\; B = 5,\; C = 3,\; D = 2$.

Expected result:
\[
  A = (10 + 5) - (3 + 2) = 15 - 5 = 10
\]

\subsection{Register Allocation}

\begin{table}[H]
  \centering
  \caption{Register assignments for Task 1}
  \renewcommand{\arraystretch}{1.3}
  \begin{tabular}{|l|l|}
    \hline
    \textbf{Register} & \textbf{Role} \\
    \hline
    \code{RAX} & Holds $A$; later holds $(A+B)$; final result \\
    \code{RBX} & Holds $B$ \\
    \code{RCX} & Holds $C$; later holds $(C+D)$ \\
    \code{R8}  & Holds $D$ (RDX reserved for \code{IDIV} remainder) \\
    \code{RSI} & String pointer for \code{sys\_write} \\
    \code{RDX} & Byte count for \code{sys\_write} \\
    \code{RDI} & File descriptor (\texttt{1} = stdout) \\
    \hline
  \end{tabular}
\end{table}

\subsection{Source Code}

\begin{lstlisting}[caption={task1.asm --- compute A = (A + B) - (C + D)}]
section .data
    lbl     db  "A = (A + B) - (C + D)", 10
    lbl_len equ $ - lbl
    pfx     db  "Result A = "
    pfx_len equ $ - pfx

section .bss
    dbuf    resb 24     ; scratch buffer for integer -> ASCII conversion

section .text
    global _start

_start:
    ; ---- print expression label ----
    mov  rax, 0x2000004
    mov  rdi, 1
    lea  rsi, [rel lbl]
    mov  rdx, lbl_len
    syscall

    ; ---- load operands ----
    mov  rax, 10        ; A = 10
    mov  rbx, 5         ; B = 5
    mov  rcx, 3         ; C = 3
    mov  r8,  2         ; D = 2

    ; ---- compute (A + B) ----
    add  rax, rbx       ; rax = 15

    ; ---- compute (C + D) ----
    add  rcx, r8        ; rcx = 5

    ; ---- A = (A+B) - (C+D) ----
    sub  rax, rcx       ; rax = 10

    ; ---- print result ----
    push rax
    mov  rax, 0x2000004
    mov  rdi, 1
    lea  rsi, [rel pfx]
    mov  rdx, pfx_len
    syscall
    pop  rax
    call print_int

    mov  rax, 0x2000001
    xor  rdi, rdi
    syscall
\end{lstlisting}

\noindent\textbf{Expected output:}
\begin{verbatim}
A = (A + B) - (C + D)
Result A = 10
\end{verbatim}

\img[0.85]{img.jpg}{Task 1 --- terminal output}

\newpage

% ===================== 5. TASK 2 =====================
\section{Task 2: Arithmetic Expression Variants 13, 8, 10}

All three expressions are evaluated in a single file (\code{task2.asm}).
The \code{IDIV} instruction performs signed division; \code{CQO} is called
before each division to sign-extend RAX into RDX:RAX.

% ----------------------------------------------------------
\subsection{Variant 13: $z = \dfrac{-a + b}{c} + 2d \Bigm/ e$}

\textbf{Values used:} $a=1,\; b=5,\; c=2,\; d=3,\; e=4$.

\textbf{Step-by-step verification:}
\begin{align*}
  -a + b &= -1 + 5 = 4 \\
  \frac{-a + b}{c} &= \frac{4}{2} = 2 \\
  2d &= 2 \times 3 = 6 \\
  2 + 6 &= 8 \\
  z &= \frac{8}{e} = \frac{8}{4} = \boxed{2}
\end{align*}

\textbf{Register trace:}

\begin{table}[H]
  \centering
  \caption{Register values during Variant 13 computation}
  \renewcommand{\arraystretch}{1.3}
  \begin{tabular}{|l|l|l|}
    \hline
    \textbf{Instruction sequence} & \textbf{RAX} & \textbf{Note} \\
    \hline
    \code{mov rax, -1; add rax, 5}  & 4    & $-a + b$ \\
    \code{cqo; idiv c (=2)}          & 2    & $(-a+b)/c$ \\
    \code{mov r8, 2; imul r8, 3}     & 2    & $2d$ in R8 \\
    \code{add rax, r8}               & 8    & $(-a+b)/c + 2d$ \\
    \code{cqo; idiv e (=4)}          & 2    & final $z$ \\
    \hline
  \end{tabular}
\end{table}

\begin{lstlisting}[caption={Variant 13 --- key assembly steps}]
    ; step 1: rax = -a + b = -1 + 5 = 4
    mov  rax, -1
    add  rax, 5

    ; step 2: rax = rax / c = 4 / 2 = 2
    cqo                     ; sign-extend RAX -> RDX:RAX
    mov  rbx, 2             ; c = 2
    idiv rbx                ; rax = 2

    ; step 3: r8 = 2 * d = 2 * 3 = 6
    mov  r8, 2
    imul r8, 3

    ; step 4: rax = 2 + 6 = 8
    add  rax, r8

    ; step 5: z = 8 / e = 8 / 4 = 2
    cqo
    mov  rbx, 4             ; e = 4
    idiv rbx                ; rax = 2
\end{lstlisting}

\newpage

% ----------------------------------------------------------
\subsection{Variant 8: $z = \dfrac{(-a+b+c+1)^3}{(a-bc+d)^2}$}

\textbf{Values used:} $a=1,\; b=2,\; c=1,\; d=3$.

\textbf{Step-by-step verification:}
\begin{align*}
  \text{numerator base} &= -a + b + c + 1 = -1 + 2 + 1 + 1 = 3 \\
  \text{numerator} &= 3^3 = 27 \\
  \text{denominator base} &= a - bc + d = 1 - 2 \cdot 1 + 3 = 2 \\
  \text{denominator} &= 2^2 = 4 \\
  z &= \frac{27}{4} = \boxed{6} \quad \text{(integer division, truncated)}
\end{align*}

\begin{lstlisting}[caption={Variant 8 --- key assembly steps}]
    ; --- numerator: (-a+b+c+1)^3 ---
    mov  rax, -1            ; -a  (a=1)
    add  rax, 2             ; -a+b  (b=2) = 1
    add  rax, 1             ; -a+b+c  (c=1) = 2
    add  rax, 1             ; -a+b+c+1 = 3
    mov  r8, rax            ; r8 = base = 3
    imul rax, r8            ; rax = 3^2 = 9
    imul rax, r8            ; rax = 3^3 = 27
    mov  r9, rax            ; r9 = numerator = 27

    ; --- denominator: (a-b*c+d)^2 ---
    mov  rax, 2             ; b = 2
    imul rax, 1             ; b*c  (c=1) = 2
    mov  rbx, 1             ; a = 1
    sub  rbx, rax           ; a - b*c = -1
    add  rbx, 3             ; a - b*c + d (d=3) = 2
    imul rbx, rbx           ; denom = 2^2 = 4

    ; --- z = numerator / denominator = 27 / 4 = 6 ---
    mov  rax, r9
    cqo
    idiv rbx                ; rax = 6
\end{lstlisting}

\newpage

% ----------------------------------------------------------
\subsection{Variant 10: $z = \dfrac{-a^2 + b^2}{a^2 - b^2 - 5}$}

\textbf{Values used:} $a=3,\; b=1$.

\textbf{Step-by-step verification:}
\begin{align*}
  a^2 &= 9, \quad b^2 = 1 \\
  \text{numerator} &= -9 + 1 = -8 \\
  \text{denominator} &= 9 - 1 - 5 = 3 \\
  z &= \frac{-8}{3} = \boxed{-2} \quad \text{(truncated toward zero)}
\end{align*}

\begin{lstlisting}[caption={Variant 10 --- key assembly steps}]
    ; compute a^2 and b^2
    mov  rax, 3             ; a = 3
    imul rax, rax           ; a^2 = 9
    mov  r8, rax            ; r8 = a^2 = 9

    mov  rax, 1             ; b = 1
    imul rax, rax           ; b^2 = 1
    mov  r9, rax            ; r9 = b^2 = 1

    ; numerator: -a^2 + b^2 = b^2 - a^2 = 1 - 9 = -8
    mov  rax, r9
    sub  rax, r8            ; rax = -8
    mov  r10, rax           ; r10 = numerator

    ; denominator: a^2 - b^2 - 5 = 9 - 1 - 5 = 3
    mov  rbx, r8            ; rbx = 9
    sub  rbx, r9            ; rbx = 8
    sub  rbx, 5             ; rbx = 3

    ; z = -8 / 3 = -2
    mov  rax, r10
    cqo                     ; sign-extend for signed division
    idiv rbx                ; rax = -2
\end{lstlisting}

\noindent\textbf{Combined terminal output of \code{task2}:}
\begin{verbatim}
--- Expr 13: z=((-a+b)/c + 2*d)/e ---
a=1 b=5 c=2 d=3 e=4
z = 2

--- Expr 8: z=(-a+b+c+1)^3 / (a-b*c+d)^2 ---
a=1 b=2 c=1 d=3
z = 6

--- Expr 10: z=(-a*a+b*b)/(a*a-b*b-5) ---
a=3 b=1
z = -2
\end{verbatim}

\img[0.85]{img2.jpg}{Task 2 --- all three expression variants}

\newpage

% ===================== 6. BCD TASK 3 =====================
\section{Task 3: BCD Conversions and Manual Arithmetic}

\subsection{Decimal to 8421 BCD Conversion}

Student number 25 is odd, so the odd column from the lab table is used:
(a) 10, (b) 13, (c) 18, (g) 44, (h) 57, (i) 69.
Each decimal digit is replaced by its 4-bit BCD code; two digits pack into one byte.

\begin{table}[H]
  \centering
  \caption{Decimal to packed 8421 BCD conversion --- odd column (student \#25)}
  \renewcommand{\arraystretch}{1.3}
  \begin{tabular}{|c|c|c|}
    \hline
    \textbf{Decimal} & \textbf{BCD (nibble notation)} & \textbf{Hex byte} \\
    \hline
    10 & \texttt{0001\ 0000} & \texttt{0x10} \\
    13 & \texttt{0001\ 0011} & \texttt{0x13} \\
    18 & \texttt{0001\ 1000} & \texttt{0x18} \\
    44 & \texttt{0100\ 0100} & \texttt{0x44} \\
    57 & \texttt{0101\ 0111} & \texttt{0x57} \\
    69 & \texttt{0110\ 1001} & \texttt{0x69} \\
    \hline
  \end{tabular}
\end{table}

\subsection{Binary vs.\ BCD Bit Comparison}

Converting the same numbers to straight binary and comparing required bits:

\begin{table}[H]
  \centering
  \caption{Binary representation vs.\ packed BCD bit width --- odd column}
  \renewcommand{\arraystretch}{1.3}
  \begin{tabular}{|c|c|c|c|c|}
    \hline
    \textbf{Decimal} & \textbf{Binary} & \textbf{Binary bits} & \textbf{Packed BCD} & \textbf{BCD bits} \\
    \hline
    10 & \texttt{0000\ 1010} & 4 & \texttt{0001\ 0000} & 8 \\
    13 & \texttt{0000\ 1101} & 4 & \texttt{0001\ 0011} & 8 \\
    18 & \texttt{0001\ 0010} & 5 & \texttt{0001\ 1000} & 8 \\
    44 & \texttt{0010\ 1100} & 6 & \texttt{0100\ 0100} & 8 \\
    57 & \texttt{0011\ 1001} & 6 & \texttt{0101\ 0111} & 8 \\
    69 & \texttt{0100\ 0101} & 7 & \texttt{0110\ 1001} & 8 \\
    \hline
  \end{tabular}
\end{table}

\noindent Packed BCD always uses exactly 8 bits for any two-digit number, while
binary uses only as many bits as the value requires. The BCD overhead is
justified when direct digit access matters more than storage efficiency.

\subsection{BCD Additions}

\subsubsection*{Example a) \quad $421 + 975$}

\[
\begin{array}{r}
  421 \\
+ 975 \\
\hline
\end{array}
\]

BCD addition digit by digit (right to left):

\begin{itemize}
  \item \textbf{Units:} $1 + 5 = 6 \leq 9$ \quad valid BCD, no correction. Digit = \textbf{6}.
  \item \textbf{Tens:} $2 + 7 = 9 \leq 9$ \quad valid BCD, no correction. Digit = \textbf{9}.
  \item \textbf{Hundreds:} $4 + 9 = 13 > 9$ \quad add $6$:
        $\texttt{1101} + \texttt{0110} = \texttt{10011}$, write digit \textbf{3}, carry \textbf{1}.
  \item \textbf{Thousands:} $0 + 0 + 1\text{(carry)} = 1$ \quad digit = \textbf{1}.
\end{itemize}

Result: $\mathbf{1396}$ \quad BCD: \texttt{0001\ 0011\ 1001\ 0110}

\subsubsection*{Example b) \quad $183 + 195$}

\begin{itemize}
  \item \textbf{Units:} $3 + 5 = 8 \leq 9$ \quad digit = \textbf{8}, carry = 0.
  \item \textbf{Tens:} $8 + 9 = 17 > 9$ \quad add $6$: $17 + 6 = 23$, digit = \textbf{7}, carry = 1.
        (In binary: $\texttt{10001} + \texttt{0110} = \texttt{10111}$,
        lower nibble = \texttt{0111} = 7, carry = 1.)
  \item \textbf{Hundreds:} $1 + 1 + 1\text{(carry)} = 3 \leq 9$ \quad digit = \textbf{3}, carry = 0.
\end{itemize}

Result: $\mathbf{378}$ \quad BCD: \texttt{0011\ 0111\ 1000}

\subsection{BCD Subtractions (10's Complement Method)}

BCD subtraction is performed by adding the 10's complement of the subtrahend.
The 10's complement of an $n$-digit number $X$ is $10^n - X$.

\subsubsection*{a) \quad $856 - 245 = 611$}

10's complement of $245$: \quad $999 - 245 + 1 = 755$

$856 + 755 = 1611$ \quad carry = 1 (end-around carry, indicates positive result)

Discard the leading 1: result = $\mathbf{611}$

\subsubsection*{b) \quad $567 - 239 = 328$}

10's complement of $239$: \quad $999 - 239 + 1 = 761$

$567 + 761 = 1328$ \quad carry = 1

Discard leading 1: result = $\mathbf{328}$

\subsubsection*{c) \quad $974 - 632 = 342$}

10's complement of $632$: \quad $999 - 632 + 1 = 368$

$974 + 368 = 1342$ \quad carry = 1

Discard leading 1: result = $\mathbf{342}$

\subsubsection*{d) \quad $453 - 983$ (negative result)}

10's complement of $983$: \quad $999 - 983 + 1 = 17$

$453 + 017 = 470$ \quad carry = 0 (no end-around carry, indicates negative result)

When carry = 0, the result is negative. Its magnitude is the 10's complement
of $470$: \quad $1000 - 470 = 530$

Result = $\mathbf{-530}$

\newpage

% ===================== 7. TASK 4 =====================
\section{Task 4: BCD Adder in Logisim}

A single-digit BCD adder can be constructed in Logisim using:

\begin{enumerate}
  \item A standard 4-bit binary full-adder (two 4-bit inputs $A$, $B$ plus
        carry-in $C_{in}$) producing a 5-bit sum $S[4:0]$.
  \item A \textbf{correction detector}: the sum exceeds valid BCD range if
        $S \geq 10$, equivalently:
        \[ C_{out} = S_4 \;\vee\; (S_3 \cdot S_2) \;\vee\; (S_3 \cdot S_1) \]
        where $S_4$ is the carry bit and $S_3$--$S_0$ are the four sum bits.
        This logic is combinational and uses two AND gates plus one OR gate.
  \item A second 4-bit adder that adds $0000$ or $0110$ (6) to the lower four
        sum bits, controlled by the correction signal $C_{out}$.
  \item The correction carry is the output carry of the second adder, forming
        the final decimal carry-out.
\end{enumerate}

A two-digit (8-bit) packed BCD adder chains two such single-digit adders, with
the carry-out of the units adder feeding the carry-in of the tens adder.

\img[0.90]{circuit.jpg}{BCD single-digit adder circuit in Logisim}

\newpage

% ===================== 8. TASK 5 =====================
\section{Task 5: BCD Addition Assembly Program}

\subsection{Specification}

Write an Assembly program that adds two 8-bit packed BCD numbers and stores
the result (including any carry) in \emph{subsequent memory locations}.

\begin{itemize}
  \item \code{bcd\_out[0]} --- the packed BCD result byte (two decimal digits).
  \item \code{bcd\_out[1]} --- the carry byte: \texttt{0} (no hundreds digit)
        or \texttt{1} (hundreds digit present).
\end{itemize}

\subsection{Algorithm}

Because the x86-64 mode removes the \code{DAA} (Decimal Adjust after Addition)
instruction that was available in 32-bit mode, the BCD correction must be
implemented manually:

\begin{enumerate}
  \item \textbf{Pre-check units digits.} Extract the low nibbles of both
        operands and sum them. If the sum $> 9$, a units correction ($+6$)
        will be needed.
  \item \textbf{Binary addition.} Add the two bytes using 16-bit arithmetic
        so intermediate values above \texttt{0xFF} are preserved in AX.
  \item \textbf{Units (low nibble) correction.} If the pre-check flagged an
        overflow, add 6 to AX to skip the six invalid BCD states.
  \item \textbf{Tens (high nibble) correction.} If the adjusted sum exceeds
        \texttt{0x9F} (i.e., the tens digit is beyond 9), add \texttt{0x60}
        to correct it and set carry $= 1$.
\end{enumerate}

\subsection{Register Allocation}

\begin{table}[H]
  \centering
  \caption{Register assignments for the BCD addition program}
  \renewcommand{\arraystretch}{1.3}
  \begin{tabular}{|l|l|}
    \hline
    \textbf{Register} & \textbf{Role} \\
    \hline
    \code{RAX / AX / AL} & Primary 16-bit accumulator; holds binary sum and BCD-corrected result \\
    \code{RBX / BX / BL} & Holds the second BCD operand (zero-extended to 16 bits) \\
    \code{RCX / CX}      & Pre-add sum of low nibbles; used to decide low correction \\
    \code{RDX / DX}      & Scratch for nibble extraction (\code{AND}/\code{SHR}) \\
    \code{R8 / R8B}      & Carry output: \texttt{0} = no carry, \texttt{1} = carry \\
    \code{RSI}           & String pointer for \code{sys\_write} \\
    \code{RDI}           & File descriptor (\texttt{1} = stdout) \\
    \code{CL}            & Scratch in \code{print\_bcd\_byte} to preserve AL \\
    \hline
  \end{tabular}
\end{table}

\subsection{Source Code}

\begin{lstlisting}[caption={bcd\_add.asm --- full BCD addition program (macOS x86-64)}]
; bcd_add_fn: adds two 8-bit packed BCD numbers
; Input : AL = first BCD byte, BL = second BCD byte
; Output: AL = BCD result, R8B = carry (0 or 1)
; Clobbers: AX, BX, CX, DX, R8

bcd_add_fn:
    movzx eax, al           ; zero-extend to 16-bit
    movzx ebx, bl

    ; Step 1: pre-check low nibbles for units overflow
    mov cx, ax
    and cx, 0x000F          ; units digit of A
    mov dx, bx
    and dx, 0x000F          ; units digit of B
    add cx, dx              ; CX = units sum (0-18)

    ; Step 2: binary addition
    add ax, bx              ; AX = binary sum

    ; Step 3: units correction if sum > 9
    cmp cx, 9
    jle .check_tens
    add ax, 6               ; add 6 to skip invalid BCD states

.check_tens:
    ; Step 4: tens correction + carry
    xor r8, r8              ; assume carry = 0
    cmp ax, 0x009F
    jle .done
    add ax, 0x0060          ; correct tens digit
    mov r8b, 1              ; carry out
    sub ax, 0x0100          ; remove carry from AX

.done:
    ret                     ; AL = BCD result, R8B = carry
\end{lstlisting}

\subsection{Worked Example}

\subsubsection*{BCD $13 + 08$ (variant pair for student \#25 --- odd=13, even=8)}

\begin{itemize}
  \item Pre-check: units digits $3 + 8 = 11 > 9$ \quad correction needed.
  \item Binary add: $\texttt{0x13} + \texttt{0x08} = \texttt{0x1B}$.
  \item Units correction: $\texttt{0x1B} + 6 = \texttt{0x21}$.
  \item Tens correction: $\texttt{0x21} \leq \texttt{0x9F}$, no correction; carry $= 0$.
  \item Result: \code{bcd\_out[0] = 0x21}, \code{bcd\_out[1] = 0}.
  \item Decimal verification: $13 + 8 = 21$ \checkmark
\end{itemize}

\subsection{Terminal Output}

\begin{verbatim}
Decimal -> BCD conversions (odd column):
10 -> 0x10
13 -> 0x13
18 -> 0x18
44 -> 0x44
57 -> 0x57
69 -> 0x69

BCD add (13 + 08): 21  carry=0
\end{verbatim}

\img[0.85]{img3.jpg}{Task 5 --- BCD addition program output}

\newpage

% ===================== 9. DEBUGGING =====================
\section{Debugging with LLDB}

The same LLDB workflow from Lab 4 applies. Key points specific to arithmetic
and BCD programs:

\begin{lstlisting}[language=bash,caption={Assembling with debug info and launching LLDB}]
nasm -f macho64 -g -F dwarf bcd_add.asm -o bcd_add.o
ld -macos_version_min 10.7 \
   -syslibroot $(xcrun --sdk macosx --show-sdk-path) \
   -lSystem -e _start bcd_add.o -o bcd_add

lldb ./bcd_add
(lldb) breakpoint set --name bcd_add_fn
(lldb) run
(lldb) register read rax rbx rcx rdx r8   ; inspect at each step
(lldb) stepi
\end{lstlisting}

\noindent\textbf{Arithmetic debugging tips:}
\begin{itemize}
  \item After every \code{IDIV}, inspect both RAX (quotient) and RDX
        (remainder). A forgotten \code{CQO} before \code{IDIV} typically
        causes a ``floating point exception'' (divide error \#DE).
  \item Watch the \textbf{overflow flag (OF)} after \code{IMUL}: if it is set,
        the product did not fit in the destination register.
  \item For BCD debugging, use \code{register read rax} after each step and
        verify the nibbles manually against the expected BCD correction.
  \item Use \code{memory read --count 2 bcd\_out} to inspect the two
        subsequent memory locations where the BCD result and carry are stored.
\end{itemize}

\newpage

% ===================== 10. CONCLUSION =====================
\section{Conclusion}

This laboratory work covered two complementary topics in Assembly programming:
signed integer arithmetic and BCD number representation.

\textbf{Task 1} demonstrated that even a simple formula like
$A = (A + B) - (C + D)$ requires explicit register management, careful
ordering of \code{ADD} and \code{SUB}, and preservation of intermediate values
across syscalls (using \code{PUSH}/\code{POP}).

\textbf{Task 2} (variants 13, 8, 10) exposed the complexity of multi-step
signed expressions. The critical discipline was always calling \code{CQO}
immediately before \code{IDIV} to avoid a divide-error fault. Integer
truncation toward zero means the Assembly result for a negative quotient
(variant 10, $z = -2$) matches the mathematical definition of truncated
integer division.

\textbf{Task 3} highlighted the structural difference between binary and BCD:
the same decimal digit 13 needs only 4 bits in binary but 8 bits in packed
BCD. BCD's advantage is that conversion to a decimal string requires only
nibble extraction --- no repeated division by 10.

\textbf{Task 5} was the most involved: emulating the missing \code{DAA}
instruction in 64-bit mode required pre-checking the unit digits before
the add, so that inter-nibble carries could be detected correctly. The
result is stored in two consecutive bytes --- \code{bcd\_out[0]} for the
BCD digit pair and \code{bcd\_out[1]} for the carry --- directly modelling
how a hardware BCD adder propagates its carry-out to the next stage.

The overall insight from this lab is that arithmetic is never ``free'' in
Assembly: every operation must be chosen and sequenced by the programmer,
data widths must match the operation, and platform constraints (like the
removal of \code{DAA} in 64-bit mode) require workarounds that expose the
exact semantics of the hardware.

\end{document}
